%%%%%%%%%%%%%%%%%%%%%%%%%%%%%%%%%%%%%%%%%%%%%%%%%%%%%%%%%%%%%%%%%%%%%%
% Lớp 10 | Phương pháp tọa độ trong mặt phẳng (𝑂𝑥𝑦) | Tích vô hướng (theo tọa độ)
% Dạng 0H9?2-3: Đẳng thức về tích vô hướng hoặc độ dài
% Số câu: 3  (sinh tự động bởi tach_cau.py; từ nay thầy sửa trực tiếp file này)
%%%%%%%%%%%%%%%%%%%%%%%%%%%%%%%%%%%%%%%%%%%%%%%%%%%%%%%%%%%%%%%%%%%%%%

%%%%%%%%%%%%%%%% TRẮC NGHIỆM 4 PHƯƠNG ÁN %%%%%%%%%%%%%%%%

% nguồn: MaHoaNH_TracNghiem_Lop10
\begin{ex}%[0H9Y2-3]%[ma:63c14725]
	Cho  đường thẳng $d_1\colon 2x + y + 3=0$ và đường thẳng $d_2\colon x+2y+6=0$. Khẳng định nào sau đây là \textbf{đúng}?
	 \choice
		{\True $d_1$ cắt $d_2$} 
		{$d_1$  song song $d_2$}
		 {$d_1$ vuông góc $d_2$}
		{$d_1$ trùng  $d_2$}
	 \loigiai{
		Xét tỉ số $\dfrac{2}{1} \ne \dfrac{1}{2}$. Suy ra $d_1$ cắt  $d_2$.\\
		Hai đường thẳng $d_1, d_2$ lần lượt có hai véc-tơ pháp tuyến $\vec{n_1}=(2;1), \vec{n_2}=(1;2)$. Ta có $\vec{n_1}\cdot\vec{n_2}= 2\cdot1 +1\cdot2=4 \ne 0$. Do đó $d_1$ không vuông góc  với $d_2$.
	 }
\end{ex}


%%%%%%%%%%%%%%%% ĐÚNG/SAI %%%%%%%%%%%%%%%%

% nguồn: MaHoaNH_DungSai_Lop10
\begin{ex}%[0H9H2-3]%[ma:e590d839]
	Trong hệ tọa độ $Oxy$, cho các điểm $A(1;4)$, $B(-3;2)$. Gọi $M$ là trung điểm  $AB$. Các khẳng định sau đây đúng hay sai?
	 \choiceTF
		 {\True $\overrightarrow{MA}\cdot \overrightarrow{AB}=-MA\cdot AB$}
		{$\overrightarrow{MB}\cdot \overrightarrow{AB}=-MB \cdot  AB$}
		{\True $\overrightarrow{AM}\cdot \overrightarrow{AB}=AM \cdot  AB$}
		{\True $\overrightarrow{MA} \cdot \overrightarrow{MB}=-MA \cdot MB$} 
	\loigiai{ 
		Do $M$ là trung  điểm $AB$ $\Rightarrow M(-1;3)$.\\
		Ta  có $\overrightarrow{MA}=(2;1)$, $\overrightarrow{MB}=(-2;-1)$, $\overrightarrow{AB}=(-4;-2)$.\\
		Khi đó $AB=2\sqrt{5}  \Rightarrow MA=MB=\dfrac{AB}{2}=\sqrt{5}$.
		\begin{itemchoice}
			\itemch Ta có $\overrightarrow{MA}\cdot\overrightarrow{AB}=2\cdot  (-4)+1\cdot (-2)=-10$.\\
			 Mà $MA\cdot AB=\sqrt{5}\cdot 2\sqrt{5}=10$.\\
			Suy ra $\overrightarrow{MA}\cdot \overrightarrow{AB}=-MA\cdot AB$. 
			\itemch Ta có $\overrightarrow{MB}\cdot\overrightarrow{AB}=(-2)\cdot  (-4)+(-1)\cdot (-2)=10$.\\
			Mà $MB\cdot AB=\sqrt{5}\cdot 2\sqrt{5}=10$.\\ 
			Suy ra $\overrightarrow{MB}\cdot \overrightarrow{AB}=MB\cdot AB$. 
			\itemch Ta có $\overrightarrow{AM}=(-2;-1)$ nên $\overrightarrow{AM}\cdot\overrightarrow{AB}=(-2)\cdot (-4)+(-1)\cdot (-2)=10$.\\ 
			 Mà $AM\cdot AB=MA\cdot AB=\sqrt{5}\cdot 2\sqrt{5}=10$.\\
			Suy ra $\overrightarrow{AM}\cdot \overrightarrow{AB}=AM\cdot  AB$.
			\itemch Ta  có $\overrightarrow{MA}\cdot\overrightarrow{MB}=2\cdot (-2)+1\cdot (-1)=-5$.\\
			Mà $MA\cdot MB=\sqrt{5}\cdot  \sqrt{5}=5$.\\
			Suy ra  $\overrightarrow{MA}\cdot \overrightarrow{MB}=-MA\cdot MB$.
		\end{itemchoice}
	 }
\end{ex}

% nguồn: MaHoaNH_DungSai_Lop10
\begin{ex}%[0H9V2-3]%[ma:ff5c03ca]
	Trong hệ tọa độ $Oxy$, cho các điểm $A(0;3)$, $B(4;3)$, $C(2;0)$. Gọi $I$  là trung điểm của đoạn $OA$, $M(x;y)$ là một điểm tùy ý. Các khẳng định sau đây đúng hay sai?
	\choiceTF 
		 {$\overrightarrow{OI}\cdot\overrightarrow{IB}+\overrightarrow{AI}\cdot\overrightarrow{IB}\neq 0$}
		 {\True $AB^2+\overrightarrow{AB}\cdot\overrightarrow{BC}+\overrightarrow{AB}\cdot\overrightarrow{CA}=0$}
		 {$(\overrightarrow{IA}+\overrightarrow{IB})\cdot\overrightarrow{IO}=\dfrac{9}{2}$}
		{\True $\overrightarrow{MA}\cdot\overrightarrow{BC}+\overrightarrow{MB}\cdot\overrightarrow{CA}+\overrightarrow{MC}\cdot\overrightarrow{AB}=0$} 
	\loigiai{ 
		\begin{itemchoice}
			 \itemch Ta có $\overrightarrow{OI}\cdot\overrightarrow{IB}+\overrightarrow{AI}\cdot\overrightarrow{IB}=\overrightarrow{IB}\cdot(\overrightarrow{OI}+\overrightarrow{AI})=\overrightarrow{IB}\cdot\overrightarrow{0}=0$.
			\itemch Ta có $AB^2+\overrightarrow{AB}\cdot\overrightarrow{BC}+\overrightarrow{AB}\cdot\overrightarrow{CA}=\overrightarrow{AB}\cdot  (\overrightarrow{AB}+\overrightarrow{BC}\cdot\overrightarrow{CA})=\overrightarrow{AB}\cdot\overrightarrow{0}=0$.
			\itemch Do $I$ là trung điểm  của đoạn $OA$ nên $I\left(0;\dfrac{3}{2}\right)$.\\
			Ta có $\overrightarrow{IA}=\left(0;\dfrac{3}{2}\right)$,  $\overrightarrow{IB}=\left(4;\dfrac{3}{2}\right)$ và $\overrightarrow{IO}=\left(0;-\dfrac{3}{2}\right)$.\\
			Do đó $\overrightarrow{IA}+\overrightarrow{IB}=(4;3)$, suy  ra $(\overrightarrow{IA}+\overrightarrow{IB})\cdot\overrightarrow{IO}=-\dfrac{9}{2}$.
			\itemch Ta  có $\overrightarrow{MA}=(-x;3-y)$, $\overrightarrow{BC}=(-2;-3)$.\\
			Suy  ra $\overrightarrow{MA}\cdot\overrightarrow{BC}=2x+3y-9$.\\
			Tương tự  $\overrightarrow{MB}\cdot\overrightarrow{CA}=2x-3y+1$, $\overrightarrow{MC}\cdot\overrightarrow{AB}=-4x+8$.\\
			Vậy $\overrightarrow{MA}\cdot\overrightarrow{BC}+\overrightarrow{MB}\cdot\overrightarrow{CA}+\overrightarrow{MC}\cdot\overrightarrow{AB}=0$ 
		\end{itemchoice}
	 }
\end{ex}
